\chapter*{Afterword}
\phantomsection
\addcontentsline{toc}{chapter}{Afterword}
\markboth{Afterword}{}
\thispagestyle{plain}

\section{Letter from the First Reader}
\begin{myquote}
\textit{------------------Chapter One------------------}

After reading this masterpiece, I was completely blown away.

I think your writing has exactly the right feel for this kind of book. You really do have a gift for writing (I mean the whole technique of introducing a new concept and then using vivid examples to explain it)!

And your range of reading and vocabulary is way richer than I imagined---Moonlight Box, Star-Absorbing Skill, ``A Spray of Plum Blossoms,'' \emph{The Count of Monte Cristo}\ldots

You really did manage to explain technical topics in plain language. I never really understood entropy before (even though basically every book about information has to mention it---like \emph{Being Digital}---but that book barely explained it, which was super unfriendly to us humanities students). Now I actually get it. It's about disorder, but the information world makes things more concise.

And I thought the example about ``effect before cause'' was so interesting---the one about how ``you'd be killed before you even see it coming.''

Also, your thinking really has depth. There are so many fresh concepts. Honestly, one reason I don't like reading these kinds of books is that every book rehashes the same ideas, like old-fashioned romance movies on repeat. But you actually bring up a lot of new concepts (maybe because the times are moving forward).

One really innovative thing about your writing is that it reads like a textbook, with genuinely interesting discussion questions. I think that's really creative and fun.

About the part on consensus, I noticed you used ``meme.'' When I read \emph{The Selfish Gene}, they kept translating it as ``memetic unit.'' You definitely explain it more clearly---less stuffy and academic\ldots


\textit{------------------Chapter Two------------------}

You definitely have the talent of a great writer!

My favorite part of Chapter Two is the dual-perspective section (the refinery and the oil field)---the shift between perspectives explains the harm of data overload so clearly and comprehensively (and I really liked the Apple Watch passage in particular. I think the layout is beautiful, the mix of long and short sentences has great rhythm, and there's a kind of ``detachment'' as you shift from abstract principles to concrete examples).

Who was it that said they were bad at writing? I think your writing is pretty good (it even creates a dramatic, novelistic feel)!

And your knowledge of history is quite impressive. I saw that one passage where you used the Enclosure Movement as an analogy. I think you learned your history well---I had to read it three or four times before I figured out the specific parallels in the analogy (though that might be because my British history is weak). There's an inherent critique of capitalists here, right? Or maybe it's that the underlying logic is the same.

I've noticed that people in the AI industry really love allusions and historical references (everyone goes on about the Tower of Babel, Shannon, and so on). Is this because you all had to take some ``History of Information'' course before they'd give you a job?

But as a high schooler who hates probability (can't even figure out expected values, and the T8 exam is coming up), I skipped every single math problem (sorry not sorry, lol). Like, what even are those math problems\ldots

Okay, I'll stop goofing around. Anyway, I can totally understand what you're trying to say. As a reader (it's not like I want to work in the industry anyway), I think this is more than enough. This book is excellent!

\textit{------------------Chapter Seven------------------}

This compression produces a by-product: when a model is asked ``What is the capital of China?,'' it doesn't actually need to ``know'' the fact. It just needs to give a high-probability answer based on the statistical pattern of this question co-occurring with ``Beijing'' in the training data---I think I've understood how LLMs answer questions now.

``Thought experiment''---that phrase has such an American vibe! You've really turned into an all-American boy, haven't you!

Are you hoping that future LLMs will be able to incorporate memory into their reasoning process---that they'll be able to think? (But memory is a function of the human brain. Our current philosophical view is that only the human brain can truly think. We believe that pig brains, artificial intelligence, and so on cannot think.)

I get the feeling you really do want it to become human.

\hfill ---Jiang Yanzi
\end{myquote}




\newpage
\section*{One Leaf Higher}

\begin{myquote}
I don't ask you to be far ahead of everyone else. I only ask that you stand ``one leaf'' higher. If you can do that, all the sunlight and rain in the sky are yours.

\hfill ---Cao Jiashu
\end{myquote}

\vspace{1em}

Yanzi,

You wrote, ``the information world makes entropy more concise.'' That single sentence is better than my entire book.

You skipped all the math---that's fine. There's a footnote in the prologue you probably skipped too---AI can turn energy into compute, and compute into labor. In the future, energy and compute will be the most fundamental productive forces. The unit of compute is called a Token. We will need more energy; compute will become as ubiquitous as electricity, more important and more useful than gold.

\vspace{1em}

People are anxious about AI, often because they equate ``can do work'' with ``is useful.''
You know what I've been thinking about lately? Athens. You're better at history than I am, so you probably know it well. Slaves did the labor; free citizens pursued philosophy, drama, politics.

\vspace{1em}

If AI does all the work, does that make us all free citizens?
But even among Athens' free citizens, there were still hierarchies, still competition for the archon's seat, still lawsuits, still athletes pushing themselves to the point of collapse at the Olympics. Physical labor was handed off, but not a single human affair was eliminated.

\vspace{1em}

This has nothing to do with whether AI exists. Taking it to the extreme---even if AI one day replaces all of our jobs, as long as society still exists, as long as governments can still redistribute social wealth, our lives probably won't be all that different from before. We don't need to compete with AI, because it can outwork every single one of us. The people around us will still be the same family, friends, classmates, colleagues, and strangers as before.

\vspace{1em}

You asked if I want AI to become human---no.
I don't want it to become human. I want it to remain a tool, always. A sufficiently intelligent, sufficiently powerful tool.
But even if one day it truly learns to ``think'' the way humans do, it still won't replace the real sense of connection between people.

\newpage
That love-letter story---the point isn't whether AI writes it well. The point is that you don't know whether the other person loves you. You have to find out for yourself, fight for it, risk rejection, and write those clumsy sentences yourself. As it stands, LLMs are designed to be inherently sycophantic tools---mirrors that forever say ``Yes, master.'' They will never refuse you. What gives a relationship weight is precisely the uncertainty, the waiting, the misunderstandings, the tentative reaching out, the choice to stand firm and not back down, the act of putting your heart out there knowing there may be no reply.

\vspace{1em}

When you chose your academic track in your first year of high school, you said you wanted humanities. ``I don't want to work in the industry.'' Honestly, I was a little surprised, and a little disappointed.
You're three months away from the gaokao now. What school, what major---frankly, the moment you told me you wanted humanities, I knew I couldn't offer you much guidance on this.

But I think that's perfectly fine. Follow your own instincts. A ruler is not for measuring the dragon---it's for summoning it. Whatever standard you use to measure yourself is what you'll grow into.


The future needs people who are different. The memes you and I share, those free-spirited lines of verse, those ``sorry not sorry, lol'' moments, everything that makes you \emph{you} and not someone else---AI can't learn that, and it can't copy it. Just find a place where you're willing to go all in. You have your own Token. Attention is your Token. Where you spend it matters more than how much you have. You don't have to win. Just don't surrender.

\vspace{1em}

Four parts, six spells, dozens of concepts. While writing, I wanted to give you everything I know. Every single one felt important; I couldn't bear to cut any of them. But if this book is ever truly useful to you, you'll still need to find the one or two sentences that matter most to you---to find your own ``fleeting octant.''
If I could leave only one signal, it would be this---AI can do everything within the distribution. What it's best at is being average.

\vspace{1em}

\textbf{What you can give this world lies outside the distribution.}

One leaf higher is enough---your own leaf.

\vspace{2em}
\begin{flushright}
Dad
\end{flushright}
